% Hints to exercises for VI/04.
\begin{enumerate}[label=\textbf{E\arabic*.},leftmargin=*,itemsep=0.6em]
\item For \cref{ex:vi04-target-test}, translate ``lies in \(V(I)\)'' point by point.
\item For \cref{ex:vi04-pullback}, check addition, multiplication, and constants after composition.
\item For \cref{ex:vi04-coordinates}, evaluate the target coordinate \(\bar y_i\) at \(\varphi(x)\).
\item For \cref{ex:vi04-contravariance}, apply both sides to an arbitrary function on \(Z\).
\item For \cref{ex:vi04-reconstruct}, set \(f_i=\Phi(\bar y_i)\) and use the target relations.
\item For \cref{ex:vi04-continuity}, a point belongs to the inverse image iff every pulled-back equation vanishes there.
\item For \cref{ex:vi04-principal-open}, replace ``\(f\neq0\)'' by ``\(\varphi^*f\neq0\)'' after composition.
\item For \cref{ex:vi04-parabola}, project to the first coordinate.
\item For \cref{ex:vi04-axis}, the variable \(y\) must map to zero.
\item For \cref{ex:vi04-affine-line-automorphism}, solve \(u=at+b\) for \(t\).
\item For \cref{ex:vi04-squaring}, the image is \(k[t^2]\).
\item For \cref{ex:vi04-maps-to-a1}, a map to \(\mathbb A^1\) has exactly one polynomial coordinate.
\item For \cref{ex:vi04-iso-ring}, use the affine correspondence on \((\varphi^*)^{-1}\).
\item For \cref{ex:vi04-closed-quotient}, use \(I(X)\subseteq I(Y)\).
\item For \cref{ex:vi04-product-pair}, concatenate the coordinate tuples and then project them back.
\item For \cref{ex:vi04-constant-map}, target coordinate functions all pull back to scalars.
\item For \cref{ex:vi04-graph}, use new coordinates \(y_1,\dots,y_m\) on the second factor.
\item For \cref{ex:vi04-kernel-image}, write out what \(\varphi^*(h)=0\) means on every \(x\in X\).
\end{enumerate}
